\begin{lemma}[Proved: projective partial resolutions retaining $\mathbf Q$-factorial points]
\label{D:Qfactorial-partial-resolution}
Let $S$ be a normal proper surface over an algebraically closed field.
Let $T$ be a subset of its finite singular locus, and assume that
$\mathcal O_{S,t}$ is $\mathbf Q$-factorial for every $t\in T$.
There is a proper birational morphism $q:X\to S$ which is an
isomorphism near $T$, which resolves every singular point outside
$T$, and for which $X$ is projective and $\mathbf Q$-factorial.
At each resolved point we may prescribe a projective regular
resolution which is an isomorphism on the punctured neighborhood.
\end{lemma}
\begin{proof}
Choose an affine neighborhood of each singular point which contains
no other singular point.  A prescribed projective resolution of
the local germ spreads to such a neighborhood by finite presentation;
after shrinking, properness, regularity and the isomorphism over
the punctured neighborhood hold for this spread-out morphism.
The specified local resolutions, the
identity neighborhoods at points of $T$, and the regular locus
glue along their common regular open subsets.  Properness is local
on the target, so the resulting $q$ is proper.  The surface $X$ is
normal, proper, and regular except at its unchanged points above
$T$.  It is therefore $\mathbf Q$-factorial.

We recall why a normal proper $\mathbf Q$-factorial surface is
projective, including the characteristic-independent argument.
Choose a regular resolution $h:W\to X$.  The proper regular
surface $W$ is projective
\cite[Lemma 54.16.11, tag 0C5P]{StacksDSurfaceResolution}.
Let $H$ be an effective ample Cartier divisor on $W$ containing
no exceptional component, and put $A=h_*H$.
The divisor $A$ is effective and $\mathbf Q$-Cartier.
For every integral curve $C\subset X$, the effective
$\mathbf Q$-Cartier pullback $h^*C$ is nonzero, and the projection
formula gives
\[
 A\cdot C=H\cdot h^*C>0.
\]
Moreover $h^*A=H+F$ for an effective exceptional
$\mathbf Q$-divisor $F$.  Since $h^*A\cdot F=0$, we have
\[
 A^2=(h^*A)^2=h^*A\cdot H=H^2+F\cdot H>0.
\]
The Nakai--Moishezon criterion makes a positive Cartier multiple
of $A$ ample.  This is the argument of
\cite[Lemma 2.2, p. 344]{C_Fujino2012}, also applied explicitly
in positive characteristic in
\cite[proof of Theorem 5.3, p. 50]{C_Tanaka2014}.
It proves the required projectivity of $X$.
\end{proof}

\begin{theorem}[Proved under stated hypotheses: mixed surface boundary]
\label{D:surface-mixed-boundary}
Let $S$ be a normal proper connected surface over an algebraically
closed field $k$ of characteristic $p>0$, and let $Y\subset S$
be a finite closed subset.  Suppose that every point $y\in Y$
satisfies at least one of the following conditions:
\begin{enumerate}
\item $\mathcal O_{S,y}$ is $\mathbf Q$-factorial;
\item the germ at $y$ admits a regular resolution whose reduced
exceptional divisor is one smooth elliptic curve;
\item $S$ is log canonical at $y$.
\end{enumerate}
Then $\pi_1(S\setminus Y)$ is topologically finitely generated.
The singularities retained in $S\setminus Y$ are unrestricted.
Condition (1) includes every point admitting a resolution by a
simple normal crossings rational tree, even when that point is
not log canonical, by Proposition~\ref{D:rational-tree-Qfactorial}.
\end{theorem}
\begin{proof}
Put $U=S\setminus Y$.  We leave unchanged every omitted singular
point which is locally $\mathbf Q$-factorial.  Every other omitted
point has a regular resolution with an elliptic or cyclic
exceptional divisor: for condition (2) this is assumed, and for
condition (3) it follows from
Lemma~\ref{D:lc-exceptional-configurations}, since its rational-tree
alternative is already $\mathbf Q$-factorial.
We use those resolutions at these points, blowing up a self-node
if needed, and resolve all the singularities lying in $U$.
Lemma~\ref{D:Qfactorial-partial-resolution} gives a normal
projective $\mathbf Q$-factorial surface $q:X\to S$ with precisely
these local properties.  Set $U'=q^{-1}(U)$.

Let $V\to U$ be a connected finite \'etale cover.  It is normal
and integral.  Its pullback to $U'$ is connected: it contains a
nonempty dense open identical to a nonempty open of $V$, and a
finite \'etale cover of the integral normal $U'$ is a disjoint
union of integral components, each dominating $U'$.
It follows that $\pi_1(U')\to\pi_1(U)$ is surjective.

The normalization of $S$ in $k(V)$ is finite and \'etale over $U$.
At an omitted point, restrict to a neighborhood containing no
other point of $Y$, and apply the corresponding local result.
For an elliptic exceptional curve, all extensions of its valuation
are tame by Proposition~\ref{D:single-curve-residue-tameness}.
For a log canonical cyclic exceptional divisor, the same assertion
is Theorem~\ref{D:lc-cusp-tame}, including the optional node blowup.
These are all the codimension-one components of $X\setminus U'$.
The unchanged $\mathbf Q$-factorial omitted points have codimension
two and create no further divisorial condition.
Thus every cover pulled back from $U$ is classified by
$Q(X,X\setminus U')$, and we have a surjection
\[
 Q(X,X\setminus U')\longrightarrow\pi_1(U).
\]
The source is topologically finitely generated by the projective
case of Proposition~\ref{D:fixed-boundary-positive}.  The same is
therefore true of its quotient.
\end{proof}

\begin{corollary}[Proved: log canonical omitted points on proper surfaces]
\label{D:lc-surface-boundary}
Let $S$ be a normal proper surface over an algebraically closed
field, and let $Y\subset S$ be finite.  If $S$ is log canonical
at every point of $Y$, then every connected component of
$S\setminus Y$ has topologically finitely generated fundamental
group.  No projectivity, Gorenstein, Cartier-index, or degree
restriction on covers is required.

More generally, let $Z$ be a proper scheme of dimension at most
two and let $Y\subset Z$ have codimension at least two in every
irreducible component.  The same conclusion holds if every point
above $Y$ on the normalization of every surface component
satisfies one of the three conditions in
Theorem~\ref{D:surface-mixed-boundary}.
\end{corollary}
\begin{proof}
In positive characteristic the first assertion is the theorem.
In characteristic zero it follows from
Proposition~\ref{D:fixed-boundary-positive} without any restriction
on the singularities.

For the second assertion we first remove nilpotents and then
take the finite disjoint union of the normalizations of the
irreducible components.  On each normal surface component in
positive characteristic the preceding theorem applies to the finite
inverse image of $Y$; in characteristic zero we use the established
characteristic-zero case of
Proposition~\ref{D:fixed-boundary-positive}.
On each component of dimension at most one the inverse image
of $Y$ is empty, so its fundamental group is that of a proper
scheme and is finitely generated.  These opens form a finite
proper surjective cover of $Z\setminus Y$.  Finite proper descent
for finite \'etale covers, in the form used in
Proposition~\ref{D:fixed-boundary-positive}, then gives finite
generation on each connected component downstairs.
\end{proof}
